\section{Introduction}\label{sec:introduction}

Hyperparameter optimization (HPO) is the dominant lever between an
out-of-the-box gradient-boosted decision tree (GBDT) implementation
and a competitive deployed model
\citep{probst2019tunability,bergstra2012random}. In tabular
classification, where GBDT methods such as XGBoost
\citep{chen2016xgboost}, LightGBM \citep{ke2017lightgbm}, and CatBoost
\citep{prokhorenkova2018catboost} remain the leading family
\citep{grinsztajn2022tree}, this lever simultaneously affects three
deliverables a practitioner cares about: predictive quality, the
wall-clock and energy cost of training, and the run-to-run stability
of the resulting model. Treating HPO as a single-objective scalar
search obscures these trade-offs and forces them to be re-discovered
post hoc.

A natural alternative is to formulate HPO as a multi-objective problem
and to seek a set of compromise solutions on the Pareto frontier
\citep{deb2002nsga2,morales2022survey}. Within this view, three
methodological challenges interact: (i) deciding \emph{which}
metrics enter the objective vector and \emph{in which direction} each
should be optimized; (ii) producing a sample of high-quality candidates
under a strict evaluation budget; and (iii) selecting a final
configuration in a way that is auditable rather than ad hoc.

The dissertation underlying this article \citep{ribeiro2026dissertation}
addressed (i)--(iii) with a Design of Experiments (DoE) front end, a
PCA + Varimax factor reduction of the response space (the
\emph{Varimax Rotated Factor Score} representation, $\VRF$), a
quadratic Response Surface Methodology (RSM) surrogate, and a
weighted-sum scalarization sweep over a $\beta$-grid that the
dissertation labelled ``NBI''. The post-defense audit summarized in
\texttt{docs/METHODOLOGY\_DECISIONS.md} of the companion repository
identified that the implementation was, in fact, a normalized
weighted-sum scalarization rather than the Das--Dennis Normal Boundary
Intersection \citep{das1998nbi}. Independently, the peer-reviewed
multivariate $\NBI$ formulation of \citet{pereira2025eaai} reframes
the same building blocks (rotated factor scores; payoff matrix;
mixture-based post-optimization on the simplex of weights) and
demonstrates the value of a true sub-problem-based $\NBI$ when
objectives are correlated.

This article presents the post-dissertation methodological refinement
of that pipeline as implemented on the
\texttt{repo-publication-readiness} branch of the companion repository.
The contributions are:
\begin{itemize}[leftmargin=1.2em]
  \item \textbf{Explicit objective directions and FMSE wrapping.}
        Every objective entering the optimization carries a required
        direction (\textsc{minimize}, \textsc{maximize},
        \textsc{target}) and an optional FMSE transform that converts
        target-seeking objectives to a uniform minimization convention,
        following Eqs.~21--22 of \citet{pereira2025eaai}.
  \item \textbf{True $N$-objective NBI.} Anchors are computed by
        per-objective surrogate optimization; the payoff matrix
        $\bm{\Phi}$ and the quasi-normal direction
        $\hat{\bm n} = -\bm{\Phi}\bm{1}/\|\bm{\Phi}\bm{1}\|$ are
        first-class objects; each sub-problem
        $\max\,t\;\text{s.t.}\;\hat{\bm F}(\bm x) = \bm F^{U} +
        \bm{\Phi}\bm\beta + t\,\hat{\bm n}$ is solved with vector
        equality constraints and SLSQP. The mathematical core is
        dimension-agnostic, supporting $q\ge 2$ objectives with no
        bi-objective hard-coding.
  \item \textbf{Design-aware surrogate modeling.} Process-quadratic
        RSM is used on factorial / central composite designs;
        Scheff\'{e} canonical mixture polynomials are used on
        simplex-lattice designs and the post-optimization stage. Cross
        usage is guarded by a validator that warns or rejects
        incompatible combinations.
  \item \textbf{Conditional MBPA post-optimization.} A second stage
        triggers only when the primary frontier diagnostics indicate a
        flat or compressed front; the trigger decision is logged
        regardless. When fired, it fits Scheff\'{e} mixture
        surrogates of generalized distance to utopia and Shannon
        entropy over the simplex of weights and refines $\bm w^{*}$
        inside an elliptical interior constraint.
  \item \textbf{Reproducibility and cost discipline.} A YAML config
        layer (Pydantic v2) makes every run declarative; a calibrated
        cost estimator predicts CPU-hours, wall-clock time, and cloud
        spend before committing to a campaign; a CI pipeline runs the
        $N$-objective NBI residual checks on every push.
\end{itemize}

The first article-track experiment, described in
Section~\ref{sec:experimental-design}, evaluates the framework on
twelve public tabular classification datasets, three GBDT algorithms,
and ten independent replicas, against four established baselines and a
legacy weighted-sum ablation. We deliberately defer the
82-dataset doctoral benchmark to a follow-up paper.

\paragraph{Outline.} Section~\ref{sec:related-work} surveys the
relevant HPO literature and positions this work against
\citet{pereira2025eaai} and \citet{ribeiro2026dissertation}.
Section~\ref{sec:method} develops the article-track method.
Section~\ref{sec:experimental-design} states the planned protocol.
Section~\ref{sec:results} reports results
\todoblock{after the v1 campaign runs}.
Section~\ref{sec:discussion} interprets them and
Section~\ref{sec:conclusion} closes.
